% Hafta 5 — Yorumlanan dillerde aynı hatalar
% Derleme: py -3.12 tools/latex/derle.py docs/week-5/assets/h05-12-yorumlanan-diller.tex
\documentclass[tikz,border=8pt]{standalone}
\usepackage{cen429-cizim}
\begin{document}
\begin{tikzpicture}
  \node[baslik] (b) at (0,0) {Yorumlanan dillerde aynı hatalar};
  \node[altbaslik] at (b.south west) {Python, JavaScript, Ruby, PHP --- kök aynı: veri ile kodun karışması};
  \node[kotukutu=48mm, anchor=north west] (k1) at (0,-2.4) {\kb{eval / exec}};
  \node[etiket, anchor=west, text width=95mm, align=left, right=5mm of k1]
    {dizeyi KOD olarak çalıştırır --- girdiyle asla birleştirilmez};
  \node[kotukutu=48mm, below=6mm of k1] (k2) {\kb{pickle / unserialize}};
  \node[etiket, anchor=west, text width=95mm, align=left, right=5mm of k2]
    {veriyi nesneye çevirirken kod çağırır --- güvenilmez veride kullanılmaz};
  \node[kotukutu=48mm, below=6mm of k2] (k3) {\kb{os.system / shell=True}};
  \node[etiket, anchor=west, text width=95mm, align=left, right=5mm of k3]
    {kabuk ayrıştırır --- argüman dizisi kullanın};
  \node[uyarikutu=48mm, below=6mm of k3] (k4) {\kb{Düzenli ifade (ReDoS)}};
  \node[etiket, anchor=west, text width=95mm, align=left, right=5mm of k4]
    {geri izleme üstel büyür --- desen ve girdi uzunluğu sınırlanır};
  \node[fit=(k1)(k2)(k3)(k4), inner sep=0pt] (grp) {};
  \node[dip, text width=155mm, anchor=north west] at ($(grp.south west)+(0,-8mm)$)
    {Dil değişir, kural değişmez: girdiyi asla yorumlayıcıya kod olarak verme.};
\end{tikzpicture}
\end{document}
